\section{Introduction}
\label{sec:introduction}

For a positive integer $N$, write $[N]=\{1,\ldots,N\}$. A set
$A\subseteq[N]$ is \emph{square-difference-free} if
\[
 (A-A)\cap\{m^2:m\in\N,\ m\ge1\}=\varnothing.
\]
Thus the restriction concerns differences in the integers. Define
\[
 s(N)=\max\{|A|:A\subseteq[N]\text{ is square-difference-free}\}.
\]
Our main result is a fixed power saving in this extremal problem.

\begin{theorem}
\label{thm:main}
There are absolute constants $c>0$ and $C<\infty$ such that, for every
integer $N\ge1$ and every square-difference-free set $A\subseteq[N]$,
\[
 |A|\le C N^{1-c}.
\]
Equivalently, every subset of $[N]$ larger than $C N^{1-c}$ contains
distinct elements whose absolute difference is the square of a positive
integer.
\end{theorem}

The exponent obtained here is positive and absolute, but extremely small.
The argument does not yield a useful numerical value or an optimal exponent.
All finite constructions in the proof are fixed independently of $N$ and
$A$; their size affects the constants and the exponent.
